\documentclass[12pt, a4paper]{article}

\usepackage[utf8]{inputenc}
\usepackage{pmboxdraw}
\usepackage{lmodern}
\usepackage[T1]{fontenc}
\usepackage{geometry}
\usepackage{xcolor}
\usepackage{titlesec}
\usepackage{titletoc}
\usepackage{fancyhdr}
\usepackage{enumitem}
\usepackage{hyperref}
\usepackage{mdframed}
\usepackage{array}
\usepackage{booktabs}
\usepackage[activate=false]{microtype}
\usepackage{parskip}
\usepackage{amsmath, amssymb}
\usepackage{tcolorbox}
\tcbuselibrary{skins,breakable,listings}
\usepackage{fancyvrb}

\geometry{top=2.5cm,bottom=2.5cm,left=2.5cm,right=2.5cm,headheight=15pt}

\definecolor{azul42}{HTML}{1F3A8A}
\definecolor{azulclaro}{HTML}{E8EEF9}
\definecolor{cinza}{HTML}{555555}
\definecolor{cinzaclaro}{HTML}{F5F5F5}

\hypersetup{colorlinks=true,linkcolor=azul42,urlcolor=azul42,
  pdftitle={Regimento - Clube de Matematica 42 Sao Paulo}}

\pagestyle{fancy}
\fancyhf{}
\fancyhead[L]{\color{azul42}\small\textit{Clube de Matemática @ 42 São Paulo}}
\fancyhead[R]{\color{cinza}\small\textit{Regimento v1.0.0}}
\fancyfoot[C]{\color{cinza}\small\thepage}
\renewcommand{\headrulewidth}{0.4pt}
\renewcommand{\headrule}{\color{azul42}\hrule width\headwidth height\headrulewidth}
\renewcommand{\footrulewidth}{0pt}

\titleformat{\section}
  {\color{azul42}\Large\bfseries}{\thesection.}{0.6em}{}
  [\vspace{2pt}{\color{azul42}\titlerule[1.5pt]}\vspace{4pt}]
\titleformat{\subsection}{\color{azul42}\large\bfseries}{\thesubsection}{0.6em}{}
\titleformat{\subsubsection}{\color{cinza}\normalsize\bfseries}{}{0em}{}
\titlespacing{\section}{0pt}{18pt}{8pt}
\titlespacing{\subsection}{0pt}{12pt}{6pt}
\titlespacing{\subsubsection}{0pt}{10pt}{4pt}

\titlecontents{section}[1.5em]
  {\vspace{4pt}\color{azul42}\bfseries}{\contentslabel{1.5em}}{}
  {\titlerule*[6pt]{.}\contentspage}
\titlecontents{subsection}[3.5em]
  {\vspace{2pt}\small}{\contentslabel{2em}}{}
  {\titlerule*[6pt]{.}\contentspage}

\setlist[itemize]{leftmargin=1.5em,itemsep=3pt,parsep=0pt,topsep=4pt,
  label={\color{azul42}\textbullet}}
\setlist[enumerate]{leftmargin=1.8em,itemsep=3pt,parsep=0pt,topsep=4pt,
  label={\color{azul42}\arabic*.}}

\newtcolorbox{infobox}[1][]{enhanced,breakable,
  colback=azulclaro,colframe=azul42,boxrule=0.8pt,arc=4pt,
  left=8pt,right=8pt,top=6pt,bottom=6pt,#1}

\newtcolorbox{codebox}{
enhanced,breakable,
  colback=cinzaclaro,colframe=cinza!40,boxrule=0.5pt,arc=3pt,
  fontupper=\ttfamily\small,left=8pt,right=8pt,top=4pt,bottom=4pt}

\newenvironment{filetree}{%
  \VerbatimEnvironment
  \begin{tcolorbox}[
    enhanced, breakable,
    colback=black!2, colframe=black!40,
    boxrule=0.5pt, arc=3pt,
    left=6pt, right=6pt, top=6pt, bottom=6pt]%
  \begin{Verbatim}[fontsize=\small, baselinestretch=1.1]%
}{%
  \end{Verbatim}%
  \end{tcolorbox}%
}

\begin{document}
\thispagestyle{empty}

% CAPA
\begin{center}
  \vspace*{3cm}
  {\fontsize{72}{72}\selectfont\color{azul42}$\Sigma$}\\[0.8cm]
  {\fontsize{32}{36}\selectfont\color{azul42}\textbf{Clube de Matemática}}\\[0.3cm]
  {\fontsize{20}{24}\selectfont\color{cinza}@ 42 São Paulo}\\[1.2cm]
  {\color{azul42}\rule{0.6\textwidth}{1.5pt}}\\[0.8cm]
  {\large\color{cinza}Regimento do Grupo de Estudos}\\[0.4cm]
  {\normalsize\color{cinza}Versão 1.0.0 --- 2026}\\[1.5cm]
  \begin{minipage}{0.75\textwidth}
    \centering\itshape\color{cinza}\small
    Este regimento apresenta os princípios, valores e diretrizes operacionais do
    Clube de Matemática da 42 São Paulo. O clube tem como propósito criar um
    ambiente de aprendizado colaborativo onde cadetes possam desenvolver
    fundamentos matemáticos sólidos, fortalecer o raciocínio lógico e, acima
    de tudo, sentir que pertencem a uma comunidade.
  \end{minipage}
\end{center}

\clearpage
\tableofcontents
\clearpage

% 1. VALORES
\section{Valores Fundamentais}

Os valores do Clube de Matemática orientam nossas interações, decisões e a cultura do grupo.
Eles refletem nosso compromisso com um ambiente de aprendizado inclusivo, seguro e eficaz,
onde todo cadete pode crescer e contribuir.

\subsection{Apoio Mútuo e Responsabilidade Compartilhada}

Cada membro contribui para o esforço coletivo e reconhece seu papel como essencial para o
sucesso do grupo. Isso inclui não apenas comprometimento individual, mas também a disposição
de apoiar e encorajar colegas que estejam enfrentando dificuldades seja com um conceito
matemático ou com a permanência na 42.

\subsection{Comunicação Aberta e Construtiva}

Criar um espaço seguro onde os cadetes possam tirar dúvidas sem julgamento é central para
o clube. Ninguém deve sentir vergonha por não entender algo --- a matemática é uma jornada
e toda pergunta é válida. O feedback entre membros deve ser honesto, respeitoso e voltado
ao crescimento.

\subsection{Respeito às Diferenças}

Os cadetes da 42 vêm de formações muito diversas: alguns nunca viram Cálculo, outros têm
graduação em exatas. Essa diversidade é uma força. O clube valoriza diferentes formas de
pensar e aprender, criando um ambiente onde todos se sintam confortáveis para contribuir
independentemente do seu ponto de partida.

\subsection{Pertencimento e Presença}

Um dos objetivos centrais do clube é fazer com que o cadete sinta que tem um lugar na
42 --- pessoas que esperam por ele, um grupo ao qual pertence. Presença regular nas sessões
não é apenas uma questão de aprendizado, mas de fortalecer os laços que mantêm a
comunidade unida.

\subsection{Curiosidade e Persistência}

A matemática exige tempo e prática. O clube celebra o processo, não apenas o resultado.
Errar faz parte do aprendizado e encarar desafios difíceis com persistência é mais valioso
do que chegar rápido à resposta certa.

\clearpage

% 2. OBJETIVOS
\section{Objetivos}

Objetivos claros fornecem direção e propósito ao Clube de Matemática. Ao defini-los desde
o início, o grupo estabelece um entendimento comum sobre o que pretende alcançar e como o
sucesso será medido.

\subsection{Objetivo Principal}

\begin{infobox}
  \textbf{Desenvolver uma base sólida de raciocínio matemático e lógico entre os cadetes
  da 42 São Paulo, contribuindo diretamente para a permanência, o pertencimento e o
  desempenho no curso.}
\end{infobox}

Para alcançar esse objetivo, o clube atuará sobre três pilares:

\begin{enumerate}
  \item Fortalecimento de fundamentos matemáticos essenciais para a programação;
  \item Criação de um espaço de comunidade e apoio mútuo entre cadetes;
  \item Redução da taxa de abandono causada por dificuldades lógico-matemáticas.
\end{enumerate}

\subsection{Objetivos Específicos}

\subsubsection{Desenvolver Competências Matemáticas Essenciais}

O clube focará no desenvolvimento progressivo de conhecimentos matemáticos práticos e
aplicáveis ao contexto da 42. Por meio de sessões estruturadas, exercícios colaborativos e
desafios progressivos, os membros construirão uma base sólida nas áreas definidas no
Capítulo~5 deste regimento.

O sucesso será medido por:

\begin{itemize}
  \item Realização de ao menos uma sessão por semana de forma contínua;
  \item Participação ativa de ao menos 8 cadetes por sessão;
  \item Relatos positivos de membros sobre impacto na confiança e na permanência.
\end{itemize}

\subsubsection{Reduzir Barreiras de Entrada no Cursus}

Muitos cadetes abandonam a 42 nas primeiras semanas por dificuldades com lógica, raciocínio
abstrato ou matemática básica. O clube oferecerá suporte específico nesses pontos críticos,
criando uma rede de apoio para quem está travado e se sente sozinho.

\subsubsection{Criar Senso de Comunidade}

Além do conteúdo matemático, o clube funcionará como um ponto de encontro regular --- um
grupo menor na comunidade maior da 42, onde o cadete tem rostos conhecidos, pessoas
que torcem por ele e um motivo concreto para aparecer no campus.

\subsection{Objetivos de Consolidação}

\subsubsection{Produzir Proposta Pedagógica Formal}

Após os primeiros ciclos de atividade, o clube documentará sua metodologia, resultados e
impacto observado em um documento formal a ser apresentado à equipe pedagógica da 42 São
Paulo. O objetivo é institucionalizar o grupo, demonstrando de forma concreta sua
contribuição para a retenção e o desempenho dos cadetes no cursus.

\subsection{Realismo e Adaptabilidade}

Os objetivos deste regimento são ambiciosos, mas alcançáveis para um grupo motivado de
cadetes. Com sessões regulares, participação ativa e liderança compartilhada, é realista
construir uma comunidade matemática relevante dentro da 42 em 6 a 12 meses. Se necessário,
o grupo revisará periodicamente suas metas e condições de funcionamento.

\clearpage

% 3. METODO
\section{Método de Trabalho}

O Clube de Matemática combina aprendizado estruturado com prática colaborativa. Para
equilibrar teoria e aplicação, dois tipos complementares de sessões serão organizados:
\textbf{sessões de estudo} e \textbf{sessões de resolução de problemas}. Além das reuniões
regulares, o clube manterá um ambiente técnico compartilhado para comunicação, documentação
e troca de conhecimento.

\subsection{Sessões de Estudo}

\begin{infobox}
  \textbf{Frequência:} uma vez por semana \quad\textbf{Duração:} 2 horas
\end{infobox}

\subsubsection{Programa da Sessão}

\begin{itemize}
  \item \textbf{00:00 -- 00:10} | \textit{Abertura:} revisão do backlog e votação do que
    será discutido na sessão.
  \item \textbf{00:10 -- 01:40} | \textit{Exploração:} discussão e resolução colaborativa
    dos temas escolhidos. Quem trouxe um problema ou referência naturalmente contribui mais
    --- mas qualquer membro pode assumir a explicação.
  \item \textbf{01:40 -- 01:55} | \textit{Registro:} soluções e descobertas relevantes são
    documentadas no repositório.
  \item \textbf{01:55 -- 02:00} | \textit{Próxima sessão:} definição do facilitador
    rotativo para o próximo encontro.
\end{itemize}

\subsection{Facilitador Rotativo}

Para garantir que as sessões comecem e fluam bem, um membro assume o papel de
\textbf{facilitador} a cada semana. Esse papel é leve: garantir que a sessão comece no
horário, conduzir a votação do backlog e manter o tempo. O facilitador \textbf{não precisa
preparar conteúdo} --- o conteúdo vem do grupo.

A rotação é definida no final de cada sessão. Se o facilitador designado não puder
comparecer, qualquer outro membro assume voluntariamente.

\subsection{Ferramentas e Ambiente}

\begin{itemize}
  \item Um canal de comunicação no Slack (ou plataforma equivalente da 42) para coordenação
    e discussões;
  \item Um repositório compartilhado (Vogsphere ou GitHub) para soluções, anotações e
    materiais de referência.
\end{itemize}

\subsection{Registro de Soluções e Compartilhamento de Conhecimento}

Documentar soluções de problemas e descobertas matemáticas é parte importante do processo
de aprendizado. Os membros são incentivados a criar registros após resolver problemas
desafiadores ou descobrir uma técnica nova.

\clearpage

% 4. COMPROMISSO
\section{Compromisso}

Ao ingressar no Clube de Matemática, cada membro afirma seu compromisso com o sucesso
coletivo do grupo. Os membros se comprometem a:

\begin{itemize}
  \item Comparecer às sessões com regularidade;
  \item Preparar as sessões que liderarem com responsabilidade;
  \item Participar ativamente da resolução de problemas;
  \item Contribuir com a base de conhecimento do grupo;
  \item Apoiar colegas que estejam enfrentando dificuldades, dentro ou fora do contexto
    matemático.
\end{itemize}

A prática individual é um componente essencial do desenvolvimento matemático. Os membros
são incentivados a dedicar cerca de 1 a 3 horas por semana, fora das sessões agendadas:

\begin{itemize}
  \item Resolver problemas em plataformas como Khan Academy, Brilliant, CSES Problem Set
    ou Art of Problem Solving;
  \item Revisar soluções da sessão anterior;
  \item Explorar aplicações matemáticas no contexto dos projetos da 42.
\end{itemize}

\clearpage

% 5. AREAS
\section{Áreas de Aprendizado}

Este capítulo apresenta as áreas matemáticas centrais que guiarão as atividades do clube,
selecionadas por sua relevância direta para o cursus da 42 e para o desenvolvimento do
raciocínio lógico-matemático dos cadetes.

\subsection{Área 1 --- Lógica e Raciocínio Matemático}

A lógica formal é a espinha dorsal da programação e da matemática. Esta área cobre
proposições, conectivos lógicos, tabelas-verdade, quantificadores, implicação, equivalência
e demonstrações. Dominar lógica permite ao cadete entender o funcionamento interno dos
programas e construir argumentos rigorosos.

\subsection{Área 2 --- Matemática Discreta}

A matemática discreta é o campo mais diretamente aplicável à ciência da computação. Inclui
teoria ingênua dos conjuntos, relações, funções, combinatória, teoria dos grafos e indução
matemática. Esses tópicos aparecem constantemente em algoritmos, estruturas de dados e
análise de complexidade.

\subsection{Área 3 --- Álgebra e Aritmética}

Uma base sólida em álgebra é fundamental para qualquer área da matemática aplicada. Esta
área cobre manipulação algébrica, teoria dos números, aritmética modular, divisibilidade,
MDC, MMC e sistemas de equações.

\subsection{Área 4 --- Probabilidade e Estatística}

Probabilidade e estatística são cada vez mais relevantes na computação moderna,
especialmente em algoritmos aleatorizados, análise de dados e inteligência artificial.
Esta área abrange probabilidade básica, distribuições, esperança, variância e análise
descritiva.

\subsection{Área 5 --- Algoritmos e Complexidade}

Esta área conecta a matemática diretamente à programação. Abrange análise de algoritmos,
notação Big-$\mathcal{O}$, complexidade de tempo e espaço, recorrências e técnicas de
design de algoritmos (divisão e conquista, programação dinâmica, algoritmos gulosos).

\clearpage

% 6. RECURSOS
\section{Recursos Necessários}

\subsection{Membros do Grupo}

O grupo busca ter aproximadamente 8 a 15 membros ativos dispostos a:

\begin{itemize}
  \item Comparecer às sessões com regularidade;
  \item Contribuir para o compartilhamento de conhecimento e documentação;
  \item Participar de exercícios práticos;
  \item Colaborar construtivamente com outros membros.
\end{itemize}

\subsection{Espaço e Ambiente}

\begin{itemize}
  \item Computadores ou quadros para demonstrações e resolução colaborativa;
  \item Conexão estável à internet;
  \item Capacidade de projeção ou compartilhamento de tela para explicações.
\end{itemize}

\subsection{Plataformas de Aprendizado}

\begin{itemize}
  \item \textbf{Khan Academy} --- fundamentos e exercícios progressivos;
  \item \textbf{Brilliant.org} --- problemas interativos de matemática e lógica;
  \item \textbf{Art of Problem Solving (AoPS)} --- problemas voltados ao nível de competição;
  \item \textbf{CSES Problem Set} --- algoritmos com foco em programação;
  \item \textbf{Project Euler} --- problemas matemáticos resolvidos com programação.
\end{itemize}

\subsection{Ferramentas}

\begin{itemize}
  \item \textbf{Latex} --- um sistema de designer de documentos;
  \item \textbf{Markdown} --- uma linguagem de marcação de documentos HTML;
  \item \textbf{Git} --- um sistema de repositório de códigos (Como latex e markdown);
\end{itemize}

\clearpage

\appendix
% APENDICE A
\section{Estrutura do Repositório de Conhecimento}

\begin{filetree}
knowledge-base/
├── README.md
├── guidelines/
│   ├── template-solution
│   └── rules-of-the-base
├── solutions/
│   ├── logic/
│   ├── discrete-math/
│   ├── algebra/
│   ├── probability/
│   └── algorithm/
├── problems/
│   ├── logic/
│   ├── discrete-math/
│   ├── algebra/
│   ├── probability/
│   └── algorithm/
├── techniques/
├── explanations/
├── core-knowledge/
└── resources/
    ├── cheat-sheets/
    └── study-materials/
\end{filetree}

\begin{description}
  \item[\texttt{knowledge-base/}]
    Raiz do repositório. Ponto de entrada para todo o conteúdo estruturado de aprendizado.

  \item[\texttt{README.md}]
    Visão geral do projeto, instruções de uso e guia de contribuição.

  \item[\texttt{guidelines/}]
    Padrões e regras que definem como o conteúdo deve ser escrito e organizado no
    repositório.
    \begin{itemize}
      \item \texttt{template-solution} --- Modelo padrão para formatar cada entrada de
        solução.
      \item \texttt{rules-of-the-base} --- Convenções de nomenclatura, critérios de
        qualidade e padrões editoriais.
    \end{itemize}

  \item[\texttt{solutions/}]
    Soluções detalhadas passo a passo para os problemas, organizadas por área matemática
    (\texttt{logic/}, \texttt{discrete-math/}, \texttt{algebra/}, \texttt{probability/},
    \texttt{algorithm/}).

  \item[\texttt{problems/}]
    Problemas para prática, ainda sem solução ou propostos como desafio. Espelha exatamente
    a estrutura de subpastas de \texttt{solutions/}.

  \item[\texttt{techniques/}]
    Estratégias reutilizáveis de resolução de problemas e métodos matemáticos (ex.:
    indução, substituição, prova por contradição).

  \item[\texttt{explanations/}]
    Textos conceituais que explicam o \emph{porquê} de teoremas, demonstrações e ideias,
    indo além do procedimento mecânico.

  \item[\texttt{core-knowledge/}]
    Definições fundamentais, axiomas e teoremas que servem de base para o restante
    do repositório.

  \item[\texttt{resources/}]
    Materiais de apoio e documentos de referência rápida.
    \begin{itemize}
      \item \texttt{cheat-sheets/} --- Folhas de fórmulas e resumos de regras para
        consulta rápida.
      \item \texttt{study-materials/} --- Leituras complementares, links externos e textos
        de referência curados.
    \end{itemize}
\end{description}

\clearpage

% APENDICE B
\section{Template de Registro de Solução}

\begin{infobox}[title=\textbf{Informações Básicas}]
  \begin{itemize}[leftmargin=1.2em]
    \item \textbf{Nome do Problema:} \hrulefill
    \item \textbf{Plataforma / Fonte:} \hrulefill
    \item \textbf{Área:} \hspace{0.3em}
      $\square$ Lógica \hspace{0.3em}
      $\square$ Mat.\ Discreta \hspace{0.3em}
      $\square$ Álgebra \hspace{0.3em}
      $\square$ Probabilidade \hspace{0.3em}
      $\square$ Algoritmos \hspace{0.3em}
      $\square$ Outro: \underline{\hspace{1.5cm}}
    \item \textbf{Dificuldade:} \hspace{0.3em}
      $\square$ Fácil \hspace{0.3em}
      $\square$ Médio \hspace{0.3em}
      $\square$ Difícil
    \item \textbf{Data de Resolução:} \hrulefill
    \item \textbf{Resolvido Por:} \hrulefill
  \end{itemize}
\end{infobox}

\vspace{0.8em}
\subsubsection*{Descrição do Problema}
Qual é o objetivo? Que tipo de raciocínio é necessário? Que informações foram fornecidas?
\vspace{3.5em}

\subsubsection*{Análise Inicial}
Que padrões você identificou? Que abordagens foram tentadas inicialmente?
\vspace{3.5em}

\subsubsection*{Solução}
Explique o raciocínio matemático central, as técnicas utilizadas e, quando aplicável, o
código desenvolvido.
\vspace{3.5em}

\subsubsection*{Conceitos e Técnicas Utilizadas}
Liste os conceitos matemáticos, fórmulas ou técnicas utilizadas.
\vspace{3em}

\subsubsection*{Principais Aprendizados}
Que padrão matemático foi identificado? Que erro atrasou a solução?
\vspace{3em}

\subsubsection*{Referências (Opcional)}
Links para documentação, tutoriais ou recursos que ajudaram.

\clearpage

% APENDICE C
\section{Áreas de Aprendizado e Subtópicos}

\subsection*{Área 1 --- Lógica e Raciocínio Matemático}

\subsubsection*{Proposições e Conectivos Lógicos}
Proposições são afirmações que podem ser verdadeiras ou falsas. Os conectivos lógicos
(negação $\neg$, conjunção $\wedge$, disjunção $\vee$, implicação $\Rightarrow$ e
bicondicional $\Leftrightarrow$) permitem combinar proposições para formar argumentos
complexos.

\subsubsection*{Tabelas-Verdade e Equivalências}
Tabelas-verdade analisam o comportamento de expressões lógicas em todos os casos possíveis.
As Leis de De Morgan ($\neg(A \wedge B) \equiv \neg A \vee \neg B$) e outras equivalências
permitem simplificar expressões.

\subsubsection*{Quantificadores e Predicados}
Quantificadores universais ($\forall$, ``para todo'') e existenciais ($\exists$, ``existe'')
permitem expressar afirmações sobre conjuntos de objetos. A lógica de predicados forma a
base da linguagem matemática formal.

\subsubsection*{Técnicas de Demonstração}
Técnicas como demonstração direta, por contra positiva, por contradição e por indução são
ferramentas fundamentais do raciocínio matemático.

\subsection*{Área 2 --- Matemática Discreta}

\subsubsection*{Teoria Ingênua dos Conjuntos}
Operações como união ($A \cup B$), interseção ($A \cap B$), diferença e produto cartesiano
($A \times B$) aparecem constantemente em estruturas de dados e teoria de tipos.

\subsubsection*{Combinatória}
Permutações $P(n,k)$, combinações $\binom{n}{k}$ e o princípio da inclusão-exclusão
aparecem em análise de complexidade de algoritmos e contagem de estruturas.

\subsubsection*{Relações e Funções}
Funções são relações com propriedades específicas essenciais para compreender mapeamentos
em programação: injetividade (saídas distintas para entradas distintas), sobrejetividade
(todo elemento do contradomínio é atingido) e bijetividade (ambas as propriedades).

\subsubsection*{Teoria dos Grafos}
Grafos $G = (V, E)$ modelam relações entre objetos. Conceitos como caminhos, ciclos,
conectividade e árvores têm aplicações diretas em redes e algoritmos de busca.

\subsubsection*{Indução Matemática}
Técnica de demonstração para afirmações sobre números naturais, especialmente útil para
provar a correção de algoritmos recursivos.

\subsection*{Área 3 --- Álgebra e Aritmética}

\subsubsection*{Aritmética Modular}
$a \equiv b \pmod{m}$ estuda o comportamento dos números em relação a um módulo.
Fundamental em criptografia e hashing.

\subsubsection*{Álgebra Linear Básica}
Vetores, matrizes e transformações lineares são essenciais para computação gráfica,
aprendizado de máquina e física computacional.

\subsubsection*{Teoria dos Números}
Divisibilidade, números primos, $\gcd(a,b)$, $\text{lcm}(a,b)$ e o Algoritmo de Euclides
são a base de algoritmos criptográficos como RSA.

\subsection*{Área 4 --- Probabilidade e Estatística}

\subsubsection*{Probabilidade Básica}
$P(A) \in [0,1]$. Probabilidade condicional $P(A|B) = \frac{P(A \cap B)}{P(B)}$ e
independência são fundamentais para algoritmos aleatorizados.

\subsubsection*{Estatística Descritiva}
Médias, medianas, desvios padrão $\sigma$, percentis e visualizações são indispensáveis
para qualquer projeto de análise de dados.

\subsubsection*{Variáveis Aleatórias e Distribuições}
Distribuições binomial, Poisson e normal descrevem padrões comuns em dados reais. Esperança
$E[X]$ e variância $\text{Var}(X)$ são medidas centrais.

\subsection*{Área 5 --- Algoritmos e Complexidade}

\subsubsection*{Recorrências}
O Teorema Mestre resolve recorrências da forma $T(n) = aT(n/b) + f(n)$, permitindo analisar
algoritmos recursivos como Merge Sort
($T(n) = 2T(n/2) + \mathcal{O}(n) = \mathcal{O}(n \log n)$).

\subsubsection*{Análise de Algoritmos}
A notação Big-$\mathcal{O}$ permite comparar algoritmos de forma independente de hardware:
$\mathcal{O}(1) \subset \mathcal{O}(\log n) \subset \mathcal{O}(n) \subset
\mathcal{O}(n \log n) \subset \mathcal{O}(n^2) \subset \mathcal{O}(2^n)$.

\subsubsection*{Técnicas de Design de Algoritmos}
Divisão e conquista, programação dinâmica e algoritmos gulosos são paradigmas fundamentais
para resolver problemas computacionais eficientemente.

\end{document}
